\scanpage{99}
% archive foliation: 98
Pour que la section donnée $\sigma'$ de $P$ soit horizontale pour une conn. proj. $c$ convenable il f. et s. \unsure{a)} que l'ouvert $U$ des pts où $D$ est non dég. soit dense sur chaque fibre de $X$ sur $S$ et b) que la section qu'on vient de définir \inserted{sur $U$} du Torseur $\uline{\operatorname{Connproj}}$ sous $\underline{\omega}^2$ \deleted{\unsure{sur}} ``soit régulière sur tt $X$'' (i.e. se prolonge \inserted{en section régulière} sur tt $X$ laquelle sera unique). Cette condition est une condition purement locale en les pts de \deleted{\unsure{pôles}} \deleted{\unsure{dégénérescence, et signifie}} dégénérescence, et signifie qu'en termes d'un \unsure{splitting} de $P^1_{X/S}\underline{\omega}$, $\gamma$ \deleted{\unsure{il}} doit avoir une singularité polaire au plus d'ordre $1$, et la partie \deleted{polaire} \deleted{\unsure{et la}} principale polaire à l'ordre $1$ près comme section nulle au pt envisagé doit être déterminée ! (écrire que $\lambda^2+\frac{d\lambda}{dt}$ n'a pas de pôle en un pt où $\lambda$ \deleted{\unsure{peut en avoir}}\inserted{en a} un !).

Donc, le long de la section marquée $\sigma$ de $P$, il y a un champ d'éléments de section d'ordre $2$ transverses à $\sigma$, et les sections $\sigma'$ de $P$ \inserted{sur $X$} pour lesquelles il y a une connexion projective de $X$ rendant $\sigma'$ horizontale (laquelle sera unique) sont celles qui, aux pts d'intersection avec $\sigma$, ont \deleted{\unsure{un}}\inserted{comme} jet d'ordre $2$ celui du champ. [\uline{NB} des élts de contact can. d'ordre $1$ à trouver \unsure{chaque} fois qu'on prend un $\Omega$ inversible, et la section de $\mathbb{P}(P^1_{X/S}(\Omega))$ définie par la
